\section{Evaluation}
\label{sec:eval}

All numbers are from the seed-42 \census run (release mode, native
x86-64) unless stated; the manifest, metrics, and set-hashes are in the
artifact, and repeat runs produce byte-identical manifests. Each
research question scores one inversion from Table~\ref{tab:defaults}.

\subsection{RQ1 — does enforcement cost scale only with governed writes?}
Partly, and the boundary is itself a result. Within the gated arm,
clean writes touching no governed type run zero policy claims (the
target-type pre-filter) and median 1.3\,ms per write versus 2.7\,ms
for compliant governed writes, which pay claim evaluation plus verdict
signing and recording. But the gated arm's ungoverned writes are not
free relative to the control arm's (median 0.7\,ms): authority
intersection runs on every graph-scoped write \emph{by design} — GS3
guards where facts may land at all and is deliberately not
abstention-eligible. Zero-cost abstention is a property of the policy
gate, not of governance as a whole, and we report it at that scope.

\subsection{RQ2 — does the gated store end cleaner than the ungated one?}
Completely, on the planted set: the gated arm's final graph contains
\textbf{0 of 6} planted defects; the control arm — the same script,
gates off — contains \textbf{6 of 6}. Each refusal names its ground:
the SHACL episode gate names the missing provenance property; the
authority refusals name the chain, the graph, and what the chain
actually holds; the policy denial cites the policy; the post-state
probe (CEN-P2) is refused by the combined-state claim a pre-state gate
would pass; and the fabricated predicate is refused by the
closed-world vocabulary policy, which is not new machinery but an
ordinary \textsc{ask} claim in \Sig requiring every predicate on the
record type to be declared.

\subsection{RQ3 — is the audit decidable in-store, and does an external
checker agree?}
\label{sec:eval-audit}
In-store, the four passes produce exactly the planted distinctions:
the dispatch inventory reports \textbf{1 violation} (an executable
tool class, ungoverned, no reason) and \textbf{1 incompleteness} (the
same shape \emph{with} a declared reason — an acknowledged bypass, not
an unknown hole); the trace audit over a five-record window reports
\textbf{0 violations and 1 incompleteness} — the deliberately
unattributed record, counted incomplete rather than misreported or
placed at the attribution root.

Against the SARC reference checker
(\texttt{besanson/sarc-governance}), we export \Sig in its spec format
and the run's 56 decisions in its flat-trace format, twice. The
\emph{faithful} export — only the evaluations \sys actually ran —
fails the reference checker with 168 discrepancies, \textbf{all of
type coverage and none of verdict, placement, or response}: the
reference invariant expects every constraint evaluated for every
action, and \sys's target-type pre-filter is invisible to it. The
\emph{padded} export, which adds explicit not-fired records for the
non-applicable constraints, passes clean. The two checkers agree
verdict-for-verdict and disagree only on coverage \emph{semantics} —
whether abstention must be materialized to be auditable — which we
report as the finding rather than resolving by fiat.

\subsection{RQ4 — does composition never widen?}
All \textbf{7 of 7} composition probes uphold the lattice contract:
the undeclared member degrades coverage to \textsf{Partial} and the
enforcement floor refuses the composed read; the cross-chain trust
pair is refused with both chains named; the expired label reads as
absent and degrades coverage; the \texttt{no-export} obligation joins
across the set; all four clean compositions pass the floor —
\textbf{zero false refusals}; the overlay rebind is refused bind-once;
and the imported pack re-interns with its content hash stable across
id spaces.

\subsection{RQ5 — do decisions replay as of their instant?}
\label{sec:eval-replay}
The amendment supersedes the label policy's claim mid-run, so every
phase-2 decision predates a rule change. Replaying all 56 recorded
decisions: \textbf{50 of 50} satisfied verdicts re-derive faithfully
under the claim and data as of their instant — and \textbf{all 50}
evaluate \emph{unsatisfied} under the amended claim, the number that
separates ``the runtime got it wrong'' from ``the spec moved,'' and
exactly what a latest-only rule set would misreport. The
\textbf{6} denials replay differently, and honestly so: their staged
deltas were rolled back — GS2 keeps the verdict, deliberately not the
attempt — so replay verifies the rules in force at the instant
(6 of 6) rather than re-deriving the outcome. Full denial
re-derivation would require traces to carry the attempted delta,
which the governed writer's traces do and the store deliberately does
not; the asymmetry is a property of refusal-by-rollback, not a gap in
bitemporality.

\subsection{The agent arm: strictness borne by the writer}
\label{sec:eval-agent}
As an extension (single model, single task, $n{=}1$ — an existence
demonstration, not a rate), an LLM agent was given the recording task
cold and authored a five-action recording without seeing \Sig.
Attempt one: 2 accepted, 3 refused — a tally with no label, a write
into a district outside the writer's authority, a fabricated
predicate. Revising strictly from the refusal text — the only channel
the writer gets — attempt two landed \textbf{5 of 5} with zero
final-graph defects: the retry loop the thesis predicts, converging in
one revision with no human in the loop.

\subsection{In the wild}
\label{sec:eval-wild}
Replaying the recorded hank trace through the audit, with the shipped
policy catalog as \Sig, yields $T \not\entails \Sig$: \textbf{2
violations, 6 incompleteness findings} over 5 records. The violations
are real and actionable — the guard enforced a locally-configured
blast-radius rule that is not an authored policy in \Sig, and the
finding's remediation is the thesis in one line: \emph{author it in
the store so it can be audited, or stop enforcing it}. The
incompleteness findings are the honest remainder: two catalog
constraints the window never exercised, undeclared placement on the
config-file rule, and partial attribution the runtime never recorded.
The synthetic \census confirms planted defects are caught; the wild
trace shows the same machinery surfacing an unstaged enforcement gap.

\subsection{Evidence sufficiency against an external benchmark}
\label{sec:eval-demm}
DEMM-Bench~\cite{solozobov2026demmbench} asks the converse of RQ3: not
whether decisions were \emph{correct}, but whether the records a
runtime emits \emph{suffice to reconstruct} eight decision-level
properties (actor identity, principal authority, action boundary,
policy basis, decision basis, data/resource touch, lifecycle context,
verification strength) under eight controlled degradation conditions.
Its lead diagnostic is \emph{Overclaim Rate} — declaring a decision's
evidence ``sufficient'' while a required property is not
reconstructable — which operationalises the \emph{container fallacy}:
inferring audit sufficiency from the presence of a trace, ledger, or
schema. Following the benchmark's invitation to add regimes and
scorers on its published contracts, we run \sys as a self-added ninth
regime (\texttt{benchmark/demm/}): the \census run exports its 56
recorded decisions as three-plane native records — writer-side guard
trace, signed verdict fact queried back from the store, bitemporal
policy snapshot — our transforms apply the benchmark's eight
degradation semantics as content-level deletions, and a content-only
reader reconstructs the eight properties, scored against the
benchmark's own construction oracle over 64 cases. The degradations
and the adapter are ours, so the result is a claim about \sys's
record \emph{format}, not a leaderboard comparison against the
benchmark's deliberately redacted-input reference scorer.

Three findings, all deterministic. First, the container fallacy
reproduces on \sys's evidence \emph{at its ceiling}: trace-, ledger-,
and schema-presence baselines overclaim on 87.5\% of cases (the
benchmark's own corpus: 50--75\%), precisely because \sys always
emits all three planes — content-level degradation leaves every
container present, so presence carries no information at all. Second,
mechanical validity checking is enough to stop overclaim but not to
localise it: a \sys-internal validator (field completeness,
evidence-hash recomputation, executor--chain consistency, grant
scope) reaches zero overclaim by refusing every degraded record
outright, matching the benchmark's stricter baselines. Third, the
property-level reading reconstructs \textbf{512 of 512} property
cells (mean PSA 1.0, zero overclaim, zero underclaim) — including
the two slices the benchmark reports as its candidate's hardest
(conflicting identity and action boundary, PSA 0.25 there): both are
decidable from \sys's records because the guard trace names its
tool, target, graph, and principal chain, and the signed verdict
names its policy, whose claim the store serves as of the decision
instant. What a degradation deletes is detectable \emph{as} absent,
so sufficiency never has to be guessed from container presence. One
asymmetry is itself corroborated by the export: the store
deliberately does not persist the denied writer (GS2 rolls the
attempt back), so actor identity for denials lives only in the
writer-side trace plane — the same division of labour the wild-trace
audit (\S\ref{sec:eval-wild}) reported as partial attribution.
